\begin{textsample}{POS Dim 1 – vem\_ed\_24 – Score 59.00 – vem\_ed\_24/t130.txt}  \label{ex:f1_pos_053}
Artigo de \textbf{opinião} Do Teletandem aos PCIs / Cesu: colaboração para a internacionalização em \textbf{casa} Priscilla Ferro priscilla.ferro@cps.sp.gov.br No \textbf{contexto} atual do ensino \textbf{superior}, os Intercâmbios Virtuais (Virtual Exchange) têm ganhado \textbf{papel} de protagonismo (O'Dowd, 2018).
Trata-se de práticas educacionais \textbf{que} \textbf{utilizam} tecnologias de comunicação para facilitar as \textbf{interações} entre indivíduos ou grupos geograficamente separados e, dessa \textbf{forma}, \textbf{promovem} a aprendizagem intercultural e o \textbf{desenvolvimento} de \textbf{competências} \textbf{linguísticas} e tecnológicas.
Há várias abordagens e metodologias para os Intercâmbios Virtuais, duas das abordagens mais conhecidas \textbf{são} o Teletandem (Telles e Vassallo, 2006) e o Collaborative Online International Learning (COIL), como \textbf{são} mundialmente conhecidos nossos PCIs (Projetos Colaborativos Internacionais).
Ambas têm os mesmos objetivos de \textbf{expandir} fronteiras educacionais e \textbf{promover} os \textbf{intercâmbios} cultural e \textbf{linguístico} entre estudantes de diferentes \textbf{partes} do mundo.
O Teletandem, formalizado em 2006 na Universidade Estadual Paulista (Unesp), evoluiu das práticas de tandem de aprendizagem de línguas estrangeiras (Telles e Vassallo, 2006).
Estudantes de diferentes países \textbf{utilizam} mensageiros instantâneos como o Skype e se \textbf{conectam} para \textbf{aprender} a língua um do outro, \textbf{promovendo} um \textbf{aprendizado} recíproco e enriquecedor.
Com a visão de"eu ajudo você a \textbf{aprender} a minha língua e você me ajuda a \textbf{aprender} a sua", o Teletandem \textbf{é} um \textbf{ambiente} de aprendizagem de línguas estrangeiras \textbf{que} vem se \textbf{estabelecendo} com sólidas bases de pesquisas e de práticas.
Os projetos COIL, por sua vez, ampliam a colaboração para além das línguas estrangeiras, pois abraçam \textbf{disciplinas} \textbf{variadas} (Guth e Helm, 2010).
Essa prática integra a colaboração internacional nos currículos acadêmicos e \textbf{permite} \textbf{que} estudantes, de diversas áreas do \textbf{saber}, trabalhem colaborativamente.
Nesses trabalhos, \textbf{é} \textbf{possível} haver diferentes \textbf{focos}, como, por exemplo, em \textbf{competências} \textbf{linguísticas} e interculturais ou em \textbf{conteúdos} \textbf{interdisciplinares} e transdisciplinares e, \textbf{seja} qual for, sua adoção no ensino \textbf{superior} traz benefícios \textbf{significativos} para alunos e professores. \textbf{continuação} Para os alunos, a aprendizagem \textbf{prática} e contextualizada oferece a oportunidade de aplicar \textbf{conhecimentos} teóricos em \textbf{contextos} reais, em parceria com colegas de diferentes países em projetos \textbf{que} refletem situações do mundo real (De Wit, 2011).
A \textbf{interação} com diferentes culturas desenvolve \textbf{competências} globais, como \textbf{compreensão} e \textbf{empatia} intercultural, \textbf{habilidades} \textbf{essenciais} no \textbf{mercado} de trabalho globalizado.
Soft skills como comunicação, liderança, resolução de problemas e trabalho em equipe também \textbf{são} aprimoradas.
Para os professores, os benefícios incluem novas \textbf{perspectivas} e métodos de ensino, o \textbf{que} enriquece o \textbf{conteúdo} acadêmico e \textbf{torna} as aulas mais dinâmicas.
A implementação de COIL \textbf{promove} a inovação no ensino, incentiva o uso de tecnologias educacionais e de metodologias \textbf{ativas} de aprendizagem (Rubin e Guth, 2023).
Professores envolvidos em projetos COIL têm a oportunidade de \textbf{ampliar} suas redes de contatos internacionais e de participar de pesquisas \textbf{colaborativas}, visando a publicações e, eventualmente, podem ministrar aulas em clases espejo, quando o professor de uma das instituições ministra uma aula para a \textbf{turma} da instituição parceira (o \textbf{que} rende um certificado de aula internacional).
No \textbf{entanto}, há professores de disciplinas técnicas ou \textbf{profissionais} \textbf{que} se sentem limitados pela barreira da língua estrangeira.
Para esses casos, faz-se necessário o \textbf{apoio} dos professores de línguas estrangeiras aos docentes de \textbf{conteúdos} \textbf{específicos} (Belz, 2003).
Os professores de línguas auxiliam na superação dessas barreiras \textbf{linguísticas}, garantindo uma comunicação eficaz para o \textbf{desenvolvimento} do projeto, além de trabalharem nuances culturais e práticas de comunicação intercultural, \textbf{essenciais} para o sucesso de colaborações internacionais.
Professores da Fatec \textbf{que} embarcam nos PCIs têm, portanto, uma grande vantagem, já \textbf{que} a instituição \textbf{conta} com mais de 280 professores de língua inglesa e mais de 80 de língua espanhola.
Esse corpo docente pode oferecer \textbf{apoio} focado em vocabulário \textbf{específico} e acompanhamento sobre a comunicação intercultural.
A colaboração entre professores de línguas e de \textbf{conteúdos} \textbf{específicos} \textbf{resulta} em projetos mais ricos e integrados, nos quais as aprendizagens de línguas e de \textbf{disciplinas} ocorrem de \textbf{forma} harmoniosa (Thorne, 2006).
Essa sinergia \textbf{é} um dos ingredientes para o sucesso das iniciativas de internacionalização, pois enriquece a experiência \textbf{educativa} e prepara melhor os alunos para atuarem em um mundo globalizado.
A \textbf{integração} de professores de línguas e de outras \textbf{disciplinas} nos projetos COIL \textbf{representa} um avanço \textbf{significativo} na internacionalização da educação.
Ao \textbf{alinhar} os esforços desses professores, cria-se uma educação mais inclusiva, eficaz e globalmente \textbf{conectada}.
Os Intercâmbios Virtuais \textbf{são} excelentes maneiras de compor a preparação dos estudantes para atuar em um mundo interdependente e em constante \textbf{transformação} e de se \textbf{construir} um \textbf{futuro} no qual o \textbf{aprendizado} transcenda fronteiras e \textbf{promova} a \textbf{paz} e a \textbf{compreensão} global por meio da educação.
Referências BELZ, Julie A.
Linguistic Perspectives on the Development of Intercultural Competence in Telecollaboration.
Language, Learning \& Technology.
Honolulu, v.7, n.2, p. 68-99, May 2003.
Disponível em: https://scholarspace.manoa.hawaii.edu/server/api/core/bitstreams/dec7d29d-63a7-4715-823f-66ae2e8021dd/content.
Acesso em 26 jul. 2024.
DE WIT, Hans.
Internationalisation of Higher Education: Nine Misconceptions.
International Higher Education, n.64, p. 1-6, Summer 2011.
Disponível em: https://ejournals.bc.edu/index.php/ihe/article/view/8556/8321.
Acesso em: 26 jul. 2024 GUTH, Sarah; HELM, Francesca.
Telecollaboration 2.0: Language, Literacies and Intercultural Learning in the 21st Century.
Bern: Peter Lang, 2010.
O'DOWD, Robert.
From telecollaboration to virtual exchange: state-of-the-art and the role of unicollaboration in moving forward.
Journal Of Virtual Exchange, [ S.L. ], v. 1, p. 1-23, 24 abr. 2018.
Disponível em: https://journal.unicollaboration.org/article/view/35567/33147.
Acesso em: 25 jul. 2018.
RUBIN, Jon; GUTH, Sarah.
The Guide to COIL Virtual Exchange.
New York: Routledge, 2023.
TELLES, João; VASSALO, Maria Luisa.
Foreign language learning In-Tandem: Teletandem as an alternative proposal in CALLT.
The ESPecialist, São Paulo, v. 27, n.2., p. 189-212, 2006.
Disponível em: https://revistas.pucsp.br/index.php/esp/article/view/1629/1048.
Acesso em: 25 jul. 2024.
THORNE, Steven L.
Pedagogical and praxiological lessons from internet-mediated intercultural foreign language education research.
In: Internet-mediated intercultural foreign language education.
Boston: Heinle, Cengage Learning, 2006.
Disponível em: https://scholarspace.manoa.hawaii.edu/server/api/core/bitstreams/b0445d68-2405-444b-b667-4ae3bf198c4a/content.
Acesso em: 25 jul.2024.

% matched lemmas: alinhar, ambiente, ampliar, apoio, aprender, aprendizado, ativo, casa, colaborativo, competência, compreensão, conectar, conhecimento, construir, contar, contexto, conteúdo, continuação, desenvolvimento, disciplina, educativo, empatia, entanto, específico, essencial, estabelecer, expandir, foco, forma, futuro, habilidade, integração, interação, intercâmbio, interdisciplinar, linguístico, mercado, opinião, papel, parte, paz, permitir, perspectiva, possível, profissional, promover, prático, que, representar, resultar, saber, ser, significativo, superior, tornar, transformação, turma, utilizar, variar
\end{textsample}
